\documentclass[11pt]{article}
\usepackage[T1]{fontenc}
\usepackage[utf8]{inputenc}
\usepackage{lmodern,microtype}
\usepackage[a4paper,margin=28mm]{geometry}
\usepackage{amsmath,amssymb,amsthm,mathtools,booktabs,longtable,array}
\usepackage[dvipsnames]{xcolor}
\usepackage{xr-hyper}
\usepackage[colorlinks=true,linkcolor=MidnightBlue,citecolor=MidnightBlue,urlcolor=MidnightBlue]{hyperref}
\newtheorem{theorem}{Theorem}[section]
\newtheorem{proposition}[theorem]{Proposition}
\newtheorem{lemma}[theorem]{Lemma}
\newtheorem{corollary}[theorem]{Corollary}
\theoremstyle{definition}\newtheorem{definition}[theorem]{Definition}
\theoremstyle{remark}\newtheorem{remark}[theorem]{Remark}
\newtheorem{question}[theorem]{Question}
\newtheorem{conjecture}[theorem]{Conjecture}
\DeclareMathOperator{\FPdim}{FPdim}
\DeclareMathOperator{\Irr}{Irr}
\DeclareMathOperator{\Tr}{Tr}
\DeclareMathOperator{\Rep}{Rep}
\let\Vec\relax
\DeclareMathOperator{\Vec}{Vec}
\newcommand{\C}{\mathcal C}
\newcommand{\Z}{\mathcal Z}
\newcommand{\Q}{\mathbb Q}
\newcommand{\N}{\mathbb Z_{\geq0}}
\newcommand{\R}{\mathbb R}
\newcommand{\CC}{\mathbb C}
\newcommand{\ZZ}{\mathbb Z}
\newcommand{\id}{\operatorname{id}}
\allowdisplaybreaks[2]
\setlength{\emergencystretch}{2em}

\externaldocument[B-]{finite}
\hypersetup{pdftitle={Exact verification guide for the rank-four classification}}
\title{Exact verification guide\\and electronic supplement}
\date{}
\begin{document}
\maketitle
\par\medskip\noindent\textbf{Research notice.}
This is AI-generated research. The results, proofs, code and data have not yet
been independently verified by human experts or formalized in Lean. Treat them
as provisional; automated checks establish only their stated computational scope.
\par\medskip

This guide accompanies \emph{Rank-four Grothendieck rings through
multiplicity 128}. It explains why the implementation checks the
criteria proved there, including completeness of every finite search.
It is not a second collection of case-by-case mathematical proofs.

\section{Reproduction and scope}
From the package root, install the dependencies recorded in
\path{requirements.txt}, then regenerate and check the certificates:
\begin{verbatim}
python generate.py --workers 4
python verify.py --workers 4
\end{verbatim}
The generator processes the pinned rank-four census through multiplicity
$128$. It applies the uniform criteria in the main paper and writes
one proof record per input ring. The verifier reconstructs the tensors,
checks each witness, and repeats every required bounded search.
For shared signed certificates it first regenerates the complete global
signed cover and reruns its universal arithmetic tests once. Each
unsigned ring is then checked against every compatible surviving
transform. The machine-generated report \path{verification/summary.json} records
the realized, excluded and unresolved counts; an unresolved record is
never interpreted as an exclusion. The criterion totals and complete
ring-by-ring results are supplied in
\path{verification/criterion_usage.json} and
\path{verification/ledger.csv.gz}.
The legacy tag \texttt{involution\_propagation\_42} is a certificate-schema
name. Discovery and checking apply the general residual-involution
criterion to complete low-induction branches, with no equality test against
a particular multiplication tensor.
The verifier reruns the full tensor audit against the pinned source.

The input consists of $889530$ compact parameter records in
\path{data/parameters_128.txt.gz}. The associated manifest fixes the
repository commit and file digests. Reconstruction uses the two
presentations of Appendix~\ref{B-sec:census}; no approximate numerical
fusion data are introduced. The source audit checks all fusion axioms
and pairwise nonisomorphism under all six unit-fixing permutations.
The explicit map in \path{data/through32_mapping.json.gz} preserves the
labels at multiplicities at most $32$. At multiplicities $33$ through
$128$, labels follow the pinned parameter-file order within each
multiplicity. Labels organize the records and are never premises of
an exclusion.

The fusion-ring-census repository supplies the complete rank-four data
and its coverage argument through multiplicity $128$. The local parameter
snapshot and manifest fix that source at commit
\texttt{7f1b36e74df53a8a1cf8597db45d261e408d61b2}; the checker's
input audit verifies tensor reconstruction, counts, isomorphism classes,
and digests. Census completeness is supplied by the pinned source
repository; this package checks its input records rather than generating
a second rank-four census.
Runtime depends on the machine and is recorded separately from the
mathematical certificates.

The supported environment is Python 3.10 or later with SymPy 1.14.0
and python-flint 0.8.0. FLINT accelerates exact integer and polynomial
arithmetic. The full-tensor audit and signed-table accelerator need a
C++17 compiler named \texttt{g++}.
The signed-search Python reference
is also supplied for cross-checking. No proprietary algebra system,
numerical eigensolver or optimization solver is needed. The tested
environment and source digests are recorded in the verification reports.
Compiler output is temporary and is not a trusted stored executable.

Generation can be resumed from its arithmetic, direct-application and
signed-search stages. The merged file
\path{certificates/classification.jsonl.gz} is a compressed stream with
a schema header followed by one record for each census ring. This
format permits verification without keeping all tensors or proof trees
in memory. Intermediate stage files are caches; the final verifier
checks their mathematical assertions afresh.

The separate command \texttt{python verify\_companion.py} checks the
bounded arithmetic supporting the companion paper. Those results are
not premises of the finite-classification verifier.

\section{Files and mathematical interface}
\begin{center}\small
\begin{tabular}{p{.40\linewidth}p{.53\linewidth}}
\toprule
File or module & Role\\
\midrule
\path{data/parameters_128.txt.gz} & All $889530$ parameter records, with exact tensor reconstruction.\\
\path{data/census_128_manifest.json} & Pinned source, digests, counts and tensor-audit results.\\
\path{data/through32_mapping.json.gz} & Explicit isomorphisms preserving the labels through $32$.\\
\texttt{data/models.json} & Fifteen full tensors with names and realization references.\\
\path{certificates/classification.jsonl.gz} & Streaming proof records for every input ring.\\
\path{software/census_data.py} & Pinned input, tensor reconstruction, labels and source checks.\\
\texttt{software/audit.py} & Tensor identity, axioms and exact multiplicities.\\
\texttt{software/applications.py} & Uniform applications and the full-tensor check of Larson's hypotheses.\\
\texttt{software/rings.py} & Ring recognizers, Casimirs and separating character coordinates.\\
\texttt{software/exact.py} & Rational polynomials, norm primes, field primes, exact roots and real signs.\\
\texttt{software/induction.py} & Exhaustive low restriction columns, Gram and twist-moment witnesses; quadratic-family certificates.\\
\texttt{software/signed.py} & Complete local signed searches and queries of the verified global cover.\\
\path{data/signed_universal_128.jsonl.gz} & Shared signed-orbit representatives surviving the universal criteria.\\
\path{software/global_signed.py} & Global-cover generation, exact arithmetic filters and complete replay.\\
\texttt{software/proofs.py} & Spherical reduction, global moments, signed parity and central-sign moments.\\
\texttt{software/dispatch.py} & Criterion dispatch, with no catalogue-index exclusion rules.\\
\bottomrule
\end{tabular}
\end{center}
The final ledger is generated by verification, not trusted as an input.
Identifiers organize results; they never prove an exclusion. A positive record requires a based-ring isomorphism to a full
tensor in the model table. The existence of the named category is the
cited mathematical input, not a new computer construction of associators.

The main paper proves the soundness implication in
Proposition~\ref{B-prop:certificate-sound}. Its standard inputs and proved uniform
applications justify direct recognizers. Lemmas~\ref{B-lem:norm} and
\ref{B-lem:field-prime} justify the arithmetic witnesses.
Proposition~\ref{B-prop:low-certificate} justifies the low search;
Proposition~\ref{B-prop:signed}, Corollary~\ref{B-cor:sign-certificate},
Proposition~\ref{B-prop:spherical-codegree-moment} and
Corollary~\ref{B-cor:signed-parity} justify the signed branches.
Corollaries~\ref{B-cor:joint-signed-conditions},
\ref{B-cor:parity-spherical-bound} and
\ref{B-cor:parity-restricted-enumeration} justify imposing their arithmetic,
indicator and coefficient-domain constraints jointly.
Proposition~\ref{B-prop:global-signed-cover} proves completeness of
shared signed-cover queries. Every negative
certificate ends at one of these criteria.

\section{Exact character and induction calculations}
The element $H=1+b_1+kb_2+k^2b_3$, $1\leq k\leq13$, separates the four
characters for at least one $k$. The checker verifies squarefreeness of
its characteristic polynomial and all identities $N_i=q_i(N_H)$.
For each irreducible factor $p$, the character values are the reductions
of $q_i(t)$ in $\Q[t]/(p)$. Coefficients are reduced modulo $p$ before
being used as number-field elements. This last requirement matters when
the component degree is smaller than four.

The positive component is identified using the strictly positive matrix
$N_H$. Proposed roots of a codegree polynomial are supplied as rational
coefficient vectors in this component field. The checker multiplies all
linear factors and recovers the entire polynomial, including multiplicities;
it does not accept a few sample roots. To reject a field inclusion, it
checks a degree obstruction or the complete-splitting prime certificate
from the main paper. A search failing to find such a prime is not treated
as an obstruction.

For a low column $v$, the equation $\sum_i v_id_i=D/f$ becomes a rational
linear system in three unknowns after setting $v_0=1$. Row reduction
identifies free coordinates. They are enumerated through the proven bounds
\[
 0\leq v_i\leq\min\{B_{0i},\lfloor\sqrt{B_{ii}}\rfloor\}.
\]
Nonintegral or negative dependent coordinates are rejected. Two nonunit
columns may be chosen explicitly; the third is determined by the required
row sum and tested for membership in its complete column set. This is an
exact enumeration, not an integer-programming heuristic.

For each remaining low part, the checker verifies a negative entry or
principal minor of $B-LL^{\mathsf T}$, or a strictly positive excess in
one of the two-moment inequalities. Every principal minor is calculated
over $\Q$. Algebraic signs are determined by Horner bounds on a rational
isolating interval, refined until zero is separated. A zero number-field
element is recognized exactly; no tolerance replaces equality.

\section{Completeness of the signed searches}
\subsection{Self-dual tables}
Use the ten-coefficient matrices in Appendix~\ref{B-sec:census} of the
finite paper, or
construct them from the symmetric products with independent tuple
$(a,b,c,d,e,f,g,h,i,j)$. In this convention the six associativity equations
are
\begin{align*}
 ad-b^2+bg+ch-d^2-e^2+1&=0,&ae-bc+bh+ci-de-ef&=0,\\
 be-cd+dh-eg+ei-fh&=0,&af+bi-c^2+cj-e^2-f^2+1&=0,\\
 bf-ce+di-eh+ej-fi&=0,&df-e^2+gi-h^2+hj-i^2+1&=0.
\end{align*}
They are obtained by multiplying the three displayed regular matrices;
conversely their vanishing gives all commuting-matrix identities and hence
all multiplication identities. This is also tested against full matrix
associativity in the small-input cross-checks.

Every signed coefficient runs through $-N,-N+2,\ldots,N$. A basis-sign
change acts on the ten positions with binary masks
\[
 (1,2,4,1,7,1,2,4,2,4).
\]
An odd coefficient cannot vanish. Select linearly independent masks among
odd positions and require those coefficients to be positive. Every sign
orbit has a representative satisfying these conditions. Residual gauge
symmetry is harmless: the search may contain several representatives of
one orbit, but none is omitted.

The complete-domain C++ search exhausts $a,b,c,d,e,f$ first.
For $b,c\ne0$, put
\[
 H=e(d+f-a)+bc,\quad G=b^2+d^2+e^2-ad-1,
 \quad J=c^2+e^2+f^2-af-1.
\]
Set $A=e(b^2-c^2)+bc(f-d)$ and
$B=b^3e-b^2cd+b(d-f)H-ebG+ecH$.
If $A\ne0$, set $i=-B/A$; if $A=0$, require $B=0$ and exhaust $i$.
Then determine $h=(H-ci)/b$, $g=(G-ch)/b$ and $j=(J-bi)/c$.
For $b\ne0,c=0$, the fourth equation instead gives $i=J/b$;
after solving for $h,g$, exhaust every bounded $j$.
The independent Python reference still exhausts $i$ before solving
these three equations, providing a different implementation of this step.
When $b=0,c\ne0$, determine
\[
 h=\frac{d^2+e^2-ad-1}{c},\quad
 i=\frac{e(d+f-a)}c,\quad
 j=\frac{c^2+e^2+f^2-af-1}c,
\]
and use $eg=h(d-f)+ei-cd$. For $e=0$, require its right side to vanish
and exhaust $g$; otherwise determine $g$ by exact division.

Finally, when $b=c=0$, first test
\[
 ad-d^2-e^2+1=0,\quad e(a-d-f)=0,\quad af-e^2-f^2+1=0.
\]
Exhaust $h,i$ and use
$eg=h(d-f)+ei$ and $ej=eh+i(f-d)$. For $e=0$ both right sides must
vanish and \emph{both} $g,j$ are exhausted. Every completed tuple is
tested against all six equations and all original bounds and parities.
This covers each possible denominator-zero case. Both implementations retain all zero branches; the extra linear
elimination only accelerates the generic C++ branch.

\subsection{Parity domains and Diophantine elimination}
For the restricted self-dual search,
Corollary~\ref{B-cor:parity-restricted-enumeration} supplies at most eight
disjoint congruence boxes. Each box is intersected with the same
odd-coefficient positivity normalization described above. Duplicate and
empty boxes are removed. The union over all compatible indicator signs
is retained; no indicator normalization replaces the positivity gauge.
Every omitted tuple is already excluded by adjoint parity.

Within a box each coordinate has an arithmetic-progression domain of
step two or four. The implementation in
\path{software/enumerate_signed_diophantine.cpp} exhausts $b,c,d,e,f$
and, for $b,c\ne0$, uses
\[
 Ca+Ai=-B_0,
\]
with the polynomial coefficients from Appendix~\ref{B-sec:census}.
The identity holds for signed coefficients as well. Write the exact
domains as $a=a_0+s_a u$ and $i=i_0+s_i v$, where $s_a,s_i$ are
two or four. The resulting linear Diophantine equation is solved by a
greatest-common-divisor progression over the full bounded intervals
for $u,v$. A zero coefficient is handled by solving for the other
variable and exhausting the free one; if both coefficients vanish,
the entire product of the domains is retained precisely when the
constant term is zero. The remaining coordinates are then determined
by the same division identities as in the complete-domain search.

If $b\ne0,c=0$, the fourth associativity equation instead gives
$fa+bi=J_0$, treated by the same progression solver, with $j$ left
free. The branch $b=0,c\ne0$ uses the division formulas already
displayed. If $b=c=0$, then $e\ne0$ forces $a=d+f$;
for $e=0$, the first equation forces $d\ne0$ and
$a=(d^2-1)/d$. Every resulting tuple is finally checked against all
six original associativity equations and every original coefficient
domain. These eliminations discard no admissible singular branch.

An additional exact filter reduces the six associativity equations
modulo $16$ in each parity box. It enumerates all permitted residue
tuples, retaining the projection of their solutions onto the free
coordinates $(b,c,d,e,f,a,i)$. An absent outer projection or an absent
pair $(a,i)$ can be skipped before integer elimination. Every integer
solution reduces to a solution of the modular equations, so this
projection preserves completeness. Final integer equations are still
checked in full.

\subsection{Shared signed-cover queries}
Appendix~\ref{B-sec:shared-signed-cover} of the main paper proves the
completeness of a single signed cover shared among all self-dual input
rings. The source \path{software/enumerate_global_signed.cpp} enumerates
one representative of each orbit under the $48$ nonunit basis
permutations and basis-sign changes, with coefficient absolute values
at most $128$. The exact $d$-number and cyclotomicity tests exclude
impossible signed orbits before any unsigned ring is queried.
These exclusions are invariant under both actions.

The retained representatives and their Casimir determinants form
\path{data/signed_universal_128.jsonl.gz}. Its manifest,
\path{data/signed_universal_128_manifest.json}, records the bound,
the number of global orbits, the universal rejection counts, the
number of retained orbits, and source and data digests. The standalone
regeneration command is
\begin{verbatim}
python build_signed_catalog.py --workers 4 --threads 4
\end{verbatim}
The main generator creates this file when it is absent. A complete
verification recompiles the enumerator, exhausts its global cover,
reruns both universal arithmetic tests and compares the regenerated
retained list and counts exactly with the supplied data. The resulting
report is \path{verification/global_signed_catalog.json}; a digest
check alone is not substituted for this mathematical replay.

Before expanding an orbit, the query applies a necessary residue
signature test. In each of the nine adjoint-coefficient positions,
record an even coefficient modulo four and replace an odd coefficient
by one; for the remaining parameter $e$, record parity alone.
Minimize this ten-entry signature over the six nonunit permutations.
It is unchanged by basis-sign changes, and the adjoint-parity
congruences force matching orbit signatures for a compatible signed
and unsigned tensor. A mismatch therefore excludes the entire orbit.
This coarse test retains all linked indicator-sign possibilities for
the subsequent exact congruence check.

For a given unsigned ring $R$, the query expands every surviving
orbit with a compatible signature into all distinct transforms. It retains precisely those
satisfying $R$'s coefficient bounds, coefficient parities, prescribed
odd-coefficient positivity normalization, and adjoint-parity
constraints. Required determinant-prime divisibilities are checked
as well. Since the determinant is orbit-invariant, those tests may
be performed before expansion. Every retained signed table receives
the remaining character, spherical or central-moment certificate.
The cover is common to all rings, but compatibility is checked against
each entire unsigned tensor. No categorification status is inferred
from an orbit label or an unsigned label.

\subsection{Replay methods and counting conventions}
The certificate's \texttt{method} field determines its replay source
and the meaning of its count:
\begin{itemize}
\item \path{exhaustive_determinant_parity_v1} uses
\path{software/enumerate_signed.cpp}. Its \texttt{total} counts all
normalized associative tuples before the determinant and parity tests.
\item \path{parity_domains_diophantine_v2} uses
\path{software/enumerate_signed_diophantine.cpp}. Its \texttt{total}
counts parity-admissible normalized associative tuples explicitly
enumerated before the determinant tests. Tuples outside its boxes
are excluded symbolically by the congruence theorem and are not counted.
\item \path{global_signed_catalog_v3} queries the regenerated shared
cover. The field \path{retained_tables} counts compatible signed
tables left for the unsigned ring after the universal arithmetic,
coefficient-domain, determinant and adjoint-parity filters.
\end{itemize}
The first two formats record the enumerator source digest, exact
preliminary rejection counts and retained tuples. The third records
both the global source and arithmetic-filter digests, the catalog
digest, and the number of orbits surviving both the residue-signature
and required-prime filters, in the field \path{eligible_orbits}. The verifier dispatches by the recorded method and regenerates
its complete candidate set. The Python reference enumerates full
local domains and can check both the restricted union and the shared
cover on bounded examples independently.

The statistics distinguish three quantities: global signed orbits;
global orbits retained by the universal criteria; and per-ring
candidates explicitly processed by each method. Global counts occur
once in \path{verification/global_signed_catalog.json}; per-method
counts occur in \path{verification/criterion_usage.json}. Their sum
is a workload count with the stated filtering conventions, not a
count of all unrestricted associative signed tables. Overlap between
candidates compatible with different unsigned rings is expected.
Regenerating witnesses may select a different supported method and hence
change these workload counts without changing the classification.
Replaying the supplied certificate stream reproduces its recorded allocation.

\subsection{Non-self-dual tables}
Normalize duality to $[0,1,3,2]$. The six-parameter matrices are
\[
T_X=\begin{pmatrix}0&1&0&0\\1&a&b&b\\0&b&c&d\\0&b&d&c\end{pmatrix},\quad
T_Y=\begin{pmatrix}0&0&1&0\\0&b&c&d\\0&d&e&f\\1&c&e&e\end{pmatrix},\quad
T_Z=T_Y^{\mathsf T}.
\]
The four equations are
\begin{align*}
 ac-b^2+2be-c^2-d^2+1&=0,&ad-b^2+be+bf-2cd&=0,\\
 b(c-d)+d(e-f)&=0,&-c^2+d^2-e^2+f^2-1&=0.
\end{align*}
The accelerated search exhausts $c,d,e$. The last equation gives
$f^2=e^2+c^2-d^2+1$; an exact integer-square test retains both roots
allowed by the original bound and parity, counting zero once.
If $c\ne d$, the third equation determines
$b=d(f-e)/(c-d)$, subject to exact divisibility and its original
bound and parity. If $c=d$, require $d(e-f)=0$ and exhaust all
allowed $b$. The first equation determines $a$ when $c\ne0$;
the second does so when $c=0\ne d$. If both vanish, exhaust all
allowed $a$. Every completed tuple is retested against all four
equations and all original bounds and parities.
The independent Python implementation exhausts $b,c,d,e,f$ before
solving for $a$, giving a separate check of the elimination.
The two sign actions change
$(a,c,d)$ and $(b,e,f)$ independently; taking the lexicographically least
image retains exactly one representative of each such orbit.

\section{Coverage, stress tests and remaining trust}
The accepted run covers the entire $889530$-ring input census. The
numbers of signed records and explicitly enumerated candidates are
generated in \path{verification/summary.json}, with the counting
convention specified there and a breakdown by method in
\path{verification/criterion_usage.json}. The verifier also re-enumerates every
low-induction and central-sign branch. The latter checks all allowed
second indicators and all sign-object twist patterns at every real
embedding, with the same rational signs under Galois conjugation.
Indicator assignments are first intersected with all adjoint-parity
congruences. Spherical global-moment branches use the parity-refined
upper bound when an odd adjoint coefficient forces an indicator to one.
The $d$-number and cyclotomicity conditions of
Corollary~\ref{B-cor:joint-signed-conditions} are checked once globally
for shared signed-cover queries, and within the signed branch for
local searches. These steps combine
the necessary conditions before more expensive character searches.

The command \texttt{python test\_applications.py} checks the uniform
applications under basis permutations. Unspecified exclusion-classification tags and deliberately altered
application witnesses are rejected. Larson's theorem has its own explicit
certificate kind with full hypothesis checks. The extension tests also
reject false parity, global-moment and literature witnesses and compare the
Bareiss Casimir determinant with symbolic determinants.

The regression suite compares the Python and C++ signed searches on all
signed inputs through multiplicity six, compares optimized elimination
against direct small bounded enumeration, and deliberately corrupts
certificates. Omitting a signed or induction branch, altering a codegree
root, supplying a false divisibility witness or assigning a false positive
model must be rejected. The results are stored in
\texttt{verification/tests.json}.

The arithmetic relies on Python, SymPy and FLINT exact polynomial and
number-field operations, and either the supplied integer C++ enumeration
or the Python reference. The C++ accelerator guards its coefficient range so all signed
arithmetic fits in its integer type: the C++ enumerators enforce
$M\leq1000$. The census uses signed 128-bit intermediates for modular
products and 64-bit integers elsewhere. Larger signed searches require
the Python backend. Casimir determinants use fraction-free Bareiss
elimination with arbitrary-precision Python integers; every exact
division is checked. Its pivots are nonzero because
$B=\sum_iN_iN_i^{\mathsf T}$ is positive definite (the unit matrix
is one summand). These are ordinary reproducible computer-assisted proofs, not
formally certified machine code or proof-assistant formalizations. The
categorical lemmas and published realizations remain mathematical inputs
subject to normal independent review. No previously stored classification
status is among those inputs.
\end{document}
